%%%%%%%%%%%%%%%%%%%%%%%%%%%%%%%%%%%%%%%%%%%%%%%%%%%%%%%%%%%%%%%%%%%%%%%%%%%
% Anhang A: Material und Requisiten im Semester (nur Skript)
% Zusammenfassung der Vorbereitungsboxen LE1--LE13. Bei Aenderungen an
% einer Vorbereitungsbox diese Liste mitpflegen.
%%%%%%%%%%%%%%%%%%%%%%%%%%%%%%%%%%%%%%%%%%%%%%%%%%%%%%%%%%%%%%%%%%%%%%%%%%%
Diese Liste fasst die Vorbereitungsboxen der 13 Lehreinheiten zusammen.
Einzelheiten zu Mengen, Proben und Plan~B stehen in der jeweiligen
Einheit. Kommerzielle Requisiten sollten vor Semesterbeginn vorhanden sein;
die Spalte \emph{Vorlauf} nennt, wie früh mit Beschaffung oder Proben
begonnen werden sollte.

\section{Requisiten und Material je Lehreinheit}

{\small
\renewcommand{\arraystretch}{1.25}
\rowcolors{2}{rwucyan40!20}{white}
\begin{longtable}{@{}>{\bfseries\color{rwuviolet}}p{7mm}>{\raggedright\arraybackslash}p{0.25\textwidth}>{\raggedright\arraybackslash}p{0.43\textwidth}>{\raggedright\arraybackslash}p{0.18\textwidth}@{}}
\toprule
\rowcolor{rwuviolet}\color{white}LE & \color{white}\textbf{Kunststück} & \color{white}\textbf{Material} & \color{white}\textbf{Vorlauf} \\
\midrule
\endhead
1 & Ideomotor-Pendel & je Person ein Faden 40\,cm mit Ring oder Mutter, Ersatzpendel; optional Pendel-App (HTTPS, QR-Code); A5-Hefte für das Tagebuch des Staunens; A3-Blätter für die Gruppenarbeit & 1 Woche \\
2 & Das geheime Objekt & undurchsichtiger Becher; Erdnüsse in der Schale; Alternative ohne Allergierisiko (Erbse in der Schote); Allergieabfrage vorab; Karteikarten A6 & 1 Woche (Allergieabfrage) \\
3 & Kruskal Count & Text auf dem Übungsblatt (geprüft mit \texttt{tools/kruskal\_check.py}); versiegelte Vorhersage; Bleistifte; Bibliothekszugang und Beamer für die Live-Recherche, Bildschirmfotos als Plan~B & 2 Wochen (Absprache Bibliothek) \\
4 & Punktwolken, schöne Fälschung, hässliches Original & Punktwolken auf den Folien; selbst gebauter fiktiver Artikel; echter, schlicht gesetzter Fachartikel; drei SIFT-Quellen mit Links/QR-Code; WLAN; optional \enquote{Das letzte Abendmahl} (Begleitbuch, Anhang~K) & 2 Wochen (Artikel bauen und auswählen) \\
5 & Switchboard / Black Box & Switchboard (kommerziell); Ersatzbatterien; Tafel oder Flipchart für den 2-4-6-Test; Zettel A7 & bei Semesterbeginn; längste Einübung \\
6 & Kontrollierte Seilroutine & Seil-Set (kommerziell oder nach Standardliteratur) plus Ersatzset; Maßband oder Zollstock für Variante~B & 2 Wochen Proben \\
7 & Charpentier-Illusion & je Set drei gleiche Streichholzschachteln, eine mit Gewicht; ein Set pro 4--5 Personen plus Ersatz; Küchenwaage; ausgewählter Artikel für die Bias-Jagd & 1 Woche \\
8 & Monty Hall (optional Chop Cup) & drei undurchsichtige Becher und eine Kugel für die Vorführung; je Paar ein Becher-Set; optional Chop Cup (kommerziell); Millimeterpapier; aktuelle Zeitungsmeldungen & 1 Woche \\
9 & Das Patter & Pendel aus LE~1 oder Charpentier-Sets aus LE~7; beide Sprechtexte geprobt; Karteikarten A7; Transkript-Kopien & 1 Woche \\
10 & Unveränderliche Summe, Vorhersage-Trick & Umschlag mit Vorhersage (Jahreszahl des Termins prüfen!); Karten und Umschlag für den Vorhersage-Trick; Taschenrechner & 1 Woche \\
11 & Patter-Umbau & Ausgangs-Patter auf Folie; Kopien des schlechten Absatzes & wenige Tage \\
12 & CRED-Strategien & zwei Vortragsausschnitte (Rechte und Abspielbarkeit prüfen); Zeitplan für die Präsentationsrunde; Feedbackbögen & 2 Wochen (Videos auswählen) \\
13 & Boundary Work & ein bereits gezeigtes Kunststück als Reprise (z.\,B.\ Switchboard); Ethik-Checklisten & wenige Tage \\
\bottomrule
\end{longtable}}

\section{Kopiervorlagen}

Alle Handouts und Übungsblätter stehen in \texttt{handouts.pdf} (Befehl
\texttt{make handouts}); diese Fassung enthält keine Lösungs- und
Regiehinweise. Pro Einheit wird ein Handout (eine Seite) und ein
Übungsblatt (zwei bis drei Seiten) je teilnehmender Person kopiert. Für
Gruppenübungen kommen A3-Blätter, Karteikarten oder zusätzliche Kopien
einzelner Materialien hinzu; die Mengen stehen in der jeweiligen
Vorbereitungsbox.

\section{Offene Auswahlentscheidungen}

Einige Materialien wählt die Lehrperson vor dem jeweiligen Termin selbst
aus, weil sie aktuell sein oder zum Fach der Gruppe passen sollen:

\begin{checkliste}
  \item LE~3: lizenzierte Datenbanken und Zugang im Hörsaal (Bibliothek)
  \item LE~4: die drei SIFT-Quellen, die \enquote{schöne Fälschung}, das
    \enquote{hässliche Original}
  \item LE~6: optional ein echter Artikel statt des fiktiven Falls
  \item LE~7: der Artikel für die Bias-Jagd
  \item LE~8: aktuelle Zeitungsmeldungen mit Statistik; Entscheidung für
    oder gegen die Chop-Cup-Routine
  \item LE~9: Pendel oder Charpentier-Illusion als Patter-Effekt
  \item LE~12: die beiden Vortragsausschnitte
  \item LE~13: das Kunststück für die Reprise
\end{checkliste}
